\definecolor{steel}{RGB}{178,186,196}   % 钢灰（封面装饰色）
\begin{titlepage}
	\begin{tikzpicture}[remember picture, overlay]
		% ---- 1. 灰色渐变铺满整个封面（金属灰 → 深灰黑）----
		\shade[top color=gray!82!black, middle color=gray!72!black, bottom color=gray!45!black]
			(current page.south west) rectangle (current page.north east);

		% ---- 2. 细点阵（模拟网格，微弱纹理）----
		\foreach \x in {1,...,20} {
			\foreach \y in {1,...,29} {
				\fill[white, opacity=0.05] (\x,\y) circle[radius=0.018cm];
			}
		}

		% ---- 3. 主标题后方的柔光晕 ----
		\fill[white, opacity=0.05] ([yshift=-8.2cm]current page.north) circle[radius=7.0cm];
		\fill[white, opacity=0.05] ([yshift=-8.2cm]current page.north) circle[radius=4.6cm];

		% ---- 4. 右上：渐开线齿轮啮合（机械主题）----
		\begin{scope}[shift={([xshift=-6.2cm,yshift=-4.6cm]current page.north east)}]
			% 大齿轮（16 齿，齿顶圆 2.75 / 齿根圆 2.0 / 轴孔 0.6）
			\draw[white, line width=1.2pt, opacity=0.85] (0,0) circle[radius=2.75];
			\draw[white, line width=1.2pt, opacity=0.85] (0,0) circle[radius=2.0];
			\draw[white, line width=1.2pt, opacity=0.85] (0,0) circle[radius=0.6];
			\foreach \a in {0,22.5,...,337.5} {
				\draw[white, line width=1.2pt, opacity=0.85] (\a:2.0) -- (\a:2.75);
			}
			% 大齿轮键槽
			\draw[white, line width=1.0pt, opacity=0.55] (0.6,0) -- (1.35,0);
			% 小齿轮（10 齿，与大齿轮啮合）
			\begin{scope}[shift={(4.2,0)}]
				\draw[white, line width=1.0pt, opacity=0.7] (0,0) circle[radius=1.7];
				\draw[white, line width=1.0pt, opacity=0.7] (0,0) circle[radius=1.0];
				\draw[white, line width=1.0pt, opacity=0.7] (0,0) circle[radius=0.35];
				\foreach \a in {0,36,...,324} {
					\draw[white, line width=1.0pt, opacity=0.7] (\a:1.0) -- (\a:1.7);
				}
				% 小齿轮键槽
				\draw[white, line width=0.8pt, opacity=0.55] (0.35,0) -- (1.0,0);
			\end{scope}
			% 中心距参考线
			\draw[white, dashed, line width=0.7pt, opacity=0.35] (0,0) -- (4.2,0);
			\node[anchor=south east, text=white, opacity=0.8, font=\small\itshape] at (4.2,-3.5) {渐开线齿轮传动};
		\end{scope}

		% ---- 5. 左下：曲柄滑块机构简图（机械主题）----
		\begin{scope}[shift={([xshift=1.0cm,yshift=1.1cm]current page.south west)}, scale=1.35]
			% 机架（带剖面线）
			\draw[white, opacity=0.85, line width=1.5pt] (0,0) -- (5.6,0);
			\foreach \x in {0,0.45,...,5.4} {
				\draw[white, opacity=0.4, line width=0.8pt] (\x,0) -- (\x-0.18,-0.32);
			}
			% 滑块导轨（虚线）
			\draw[white, dashed, opacity=0.4, line width=0.8pt] (4.1,0.9) -- (5.6,0.9);
			% 曲柄
			\draw[white, line width=1.8pt, opacity=0.95] (0.7,0) -- (2.0,1.55);
			% 连杆
			\draw[white, line width=1.8pt, opacity=0.95] (2.0,1.55) -- (4.5,0.9);
			% 滑块
			\draw[white, line width=1.5pt, opacity=0.95] (4.45,0.45) rectangle (5.15,1.35);
			% 固定铰链符号（V 形机架线）
			\draw[white, line width=1.0pt, opacity=0.8] (0.7,0) -- (0.98,-0.28) -- (1.26,0);
			% 转动副（铰链中心）
			\fill[white] (0.7,0) circle[radius=0.12];
			\fill[white] (2.0,1.55) circle[radius=0.12];
			\fill[white] (4.5,0.9) circle[radius=0.12];
			% 标签
			\node[anchor=north west, text=white, opacity=0.8, font=\small\itshape] at (0,2.9) {曲柄滑块机构};
		\end{scope}

		% ---- 6. 银色细边框 ----
		\draw[steel, line width=0.7pt, opacity=0.5]
			([shift={(0.85,0.85)}]current page.south west) rectangle
			([shift={(-0.85,-0.85)}]current page.north east);

		% ---- 7. 顶部眉题（英文小字 + 银色短线）----
		\coordinate (T) at ([yshift=-3.5cm]current page.north);
		\draw[steel, line width=1.0pt] (T) ++(-3.4cm,0) -- ++(1.0cm,0);
		\draw[steel, line width=1.0pt] (T) ++(2.4cm,0) -- ++(1.0cm,0);
		\node[anchor=north, text=white, opacity=0.85, align=center] at (T) {
			{\sffamily\fontsize{13}{16}\selectfont
				\uppercase{Course \ Notes}}
		};

		% ---- 8. 主标题（居中，使用 \notename）----
		\node[anchor=center, align=center, text=white] at ([yshift=-8.0cm]current page.north) {%
			{\fontsize{40}{50}\sffamily\bfseries \notename}
		};

		% ---- 9. 银色菱形 + 双细线 ----
		\coordinate (C) at ([yshift=-11.4cm]current page.north);
		\draw[steel, line width=1.1pt] (C) ++(-2.2cm,0) -- ++(1.55cm,0);
		\draw[steel, line width=1.1pt] (C) ++(0.65cm,0) -- ++(1.55cm,0);
		\node[fill=steel, minimum size=0.24cm, inner sep=0pt, rotate=45] at (C) {};

		% ---- 10. 副标题 ----
		\node[anchor=north, text=white, opacity=0.9, align=center]
			at ([yshift=-12.2cm]current page.north) {
			{\fontsize{16}{20}\sffamily 学\ \ 习\ \ 笔\ \ 记}
		};

		% ---- 11. 底部作者、日期信息 ----
		\coordinate (B) at ([yshift=3.4cm]current page.south);
		\draw[steel, line width=0.8pt, opacity=0.65] (B) ++(-1.1cm,0.85cm) -- ++(2.2cm,0);
		\node[anchor=south, text=white, align=center] at (B) {
			\begin{tabular}{c}
				{\fontsize{20}{24}\sffamily\bfseries \noteauthor} \\[0.4cm]
				{\large \sffamily \today}
			\end{tabular}
		};
	\end{tikzpicture}
\end{titlepage}
